\chapter{CP-SAT: il motore che cerca, deduce, ritratta}
\label{ch:metodo_cpsat}

\section{Il problema concreto}

Immagina un dirigente scolastico la prima settimana di
settembre. Sul tavolo ha l'elenco delle classi (20--80, a
seconda della scuola), l'elenco dei docenti (40--150), i monte
ore di ogni materia in ogni indirizzo, le aule disponibili e
una pila di richieste personali (``il prof Bianchi non lavora
sabato'', ``il lab di fisica \`e occupato il marted\`i mattina
da un altro istituto''). Deve produrre un orario settimanale
in cui ogni lezione sta in uno slot compatibile, nessun
docente \`e in due posti contemporaneamente, nessuna classe
ha buchi inutili e ogni materia rispetta il monte ore.

A occhio sembra che basti ``provare''. In pratica il numero
di assegnamenti possibili \`e astronomico. Una scuola media di
30 classi con 5 ore al giorno per 6 giorni ha
$30 \times 5 \times 6 = 900$ slot da riempire scegliendo per
ognuno una lezione fra centinaia di possibili: $10^{2700}$
combinazioni circa. Nessun computer al mondo le pu\`o esaminare
una per una. Serve una macchina che, mentre prova, capisca
\emph{quali combinazioni si possono escludere subito} senza
esaminarle.

CP-SAT\index{CP-SAT} \`e questa macchina. Il nome significa
\emph{Constraint Programming over SAT}: prende l'idea
classica della programmazione a vincoli e la implementa
sopra un solver SAT moderno (cio\`e un risolutore di formule
booleane), con tutte le ottimizzazioni che la comunit\`a SAT
ha sviluppato negli ultimi vent'anni. \`E sviluppato da Google
come parte del pacchetto OR-Tools~\cite{ortools}, ed \`e oggi
uno dei migliori solver gratuiti per problemi combinatori di
medie/grandi dimensioni.

\begin{ideabox}
\textbf{In una frase}: CP-SAT \`e un motore che, dato un
elenco di variabili, di vincoli da rispettare e di un
obiettivo da minimizzare, cerca un'assegnazione di valori
alle variabili che rispetti tutto e renda l'obiettivo il pi\`u
piccolo possibile, senza enumerare le possibilit\`a una a una.
\end{ideabox}

\section{Una storia in cinque tappe}

Per capire come funziona CP-SAT \`e utile ripercorrere
brevemente cinque idee successive.

\subsection{(1) Il problema SAT puro}

Negli anni '60 i logici si chiedevano: data una formula
booleana (un'espressione di variabili true/false collegate da
\textsc{and}, \textsc{or}, \textsc{not}), esiste un assegnamento
di verit\`a che la rende vera? Questo \`e il problema SAT,
formalmente NP-completo~\cite{garey1979}. Una formula tipica
appare cos\`i:
\[
  (x_1 \vee \neg x_2 \vee x_3) \wedge (\neg x_1 \vee x_2)
  \wedge (x_2 \vee \neg x_3)
\]
Le parentesi tonde si chiamano \emph{clausole}; ogni clausola
\`e un \textsc{or} di letterali (una variabile o la sua
negazione). Una formula in questa forma \`e detta
\emph{normalizzata in CNF} (\emph{conjunctive normal form}).
Il solver deve trovare un'assegnazione che renda \emph{vera
ogni} clausola.

\subsection{(2) DPLL: backtracking + propagazione unaria}

Negli anni '60 Davis, Putnam, Logemann e Loveland proposero
l'algoritmo DPLL: scegli una variabile, prova
\texttt{x=true}, propaga le conseguenze, se trovi una
contraddizione prova \texttt{x=false}. Si chiama
\emph{backtracking}: se sbagli torni indietro. La
propagazione unaria si chiama \emph{unit propagation}: se in
una clausola tutti i letterali tranne uno sono falsi, quello
rimasto deve essere vero.

\subsection{(3) CDCL: imparare dai propri errori}

Negli anni '90 si \`e capito che il backtracking puro spreca
energia: spesso il solver torna indietro a uno stato che
\emph{contiene la stessa contraddizione} di quello appena
abbandonato, e ricomincia da capo. La svolta \`e stata il
\emph{conflict-driven clause learning}: quando trovi una
contraddizione, analizza le scelte che ti ci hanno portato e
sintetizzane una nuova clausola che \emph{vieta} quella
combinazione fatale. Aggiungi la clausola alla formula. La
prossima volta che ti avvicini a quella combinazione, la unit
propagation ti blocca subito. Questa idea, implementata in
solver come Chaff~\cite{moskewicz2001} e tutti i suoi
successori, ha portato a un salto di prestazioni di ordini di
grandezza.

\subsection{(4) Constraint Programming: variabili intere e
vincoli ad alto livello}

Nel frattempo, la comunit\`a di constraint programming
sviluppava un linguaggio diverso: invece di sole variabili
booleane e clausole, lavorava con variabili intere su domini
finiti (es. $x \in \{1,\ldots,30\}$) e vincoli ``ad alto
livello'' come \texttt{AllDifferent($x_1,\ldots,x_n$)} o
\texttt{Cumulative} per problemi di scheduling. Per ogni
vincolo si scrivono \emph{algoritmi di propagazione
specializzati} che, dato il dominio corrente di ogni
variabile, eliminano i valori che non possono pi\`u far parte
di una soluzione.

\subsection{(5) CP-SAT: la fusione}

CP-SAT prende il meglio dei due mondi:
\begin{itemize}
\item dal lato CP, eredita le variabili intere e i vincoli
  globali (il problema rimane scrivibile in modo naturale);
\item dal lato SAT, eredita CDCL, unit propagation e tutte le
  ottimizzazioni ingegneristiche (decision heuristic,
  restart, in-processing).
\end{itemize}

L'implementazione internamente traduce le variabili intere e
i vincoli ad alto livello in una rete di booleane equivalente
(``lazy clause generation''): ogni volta che il solver
prende una decisione su una variabile intera, le booleane
associate vengono propagate dal motore SAT, e quando il SAT
solver scopre una contraddizione, la nuova clausola appresa
si traduce di nuovo in vincoli sulla rete intera. \`E una
danza continua fra i due livelli.

\begin{storiabox}
La famiglia DPLL ha origine nel 1962 (Davis, Putnam,
Logemann, Loveland). CDCL nasce nel 1996--1997 con GRASP
(Marques-Silva, Sakallah). Chaff (2001) introduce le strutture
dati moderne (two-watched literals). CP-SAT come prodotto
Google nasce intorno al 2010 con OR-Tools, e ha vinto le
edizioni 2018, 2019, 2020 e 2021 della Mini\-Zinc Challenge,
la principale competizione di constraint solver.
\end{storiabox}

\section{Costruzione passo-passo: cosa serve sapere}

\subsection{Variabile in CP}

\begin{defbox}
Una \textbf{variabile} in CP \`e un'incognita con un
\emph{dominio} finito di valori possibili. Esempi:
\begin{itemize}
\item $\texttt{slot}_{\text{matematica-1A}} \in
  \{1, 2, \ldots, 36\}$ (lo slot in cui sta la prima ora di
  matematica della 1A; numero da 1 a 36 se ci sono 6 giorni
  $\times$ 6 ore).
\item $\texttt{aula}_{\text{fisica-3B-marted\`i-9}} \in
  \{\text{LabFisica1}, \text{LabFisica2}\}$.
\item $\texttt{is\_present}_{\text{Rossi-luned\`i-1}} \in
  \{0, 1\}$ (variabile booleana).
\end{itemize}
\end{defbox}

In Python con OR-Tools si scrivono cos\`i:
\begin{lstlisting}[style=py]
from ortools.sat.python import cp_model
m = cp_model.CpModel()

slot = m.NewIntVar(1, 36, "slot_mate_1A")
aula = m.NewIntVarFromDomain(
    cp_model.Domain.FromValues([101, 102]), "aula_fisica_3B_mar_9")
present = m.NewBoolVar("Rossi_lun_1")
\end{lstlisting}

\subsection{Vincolo}

\begin{defbox}
Un \textbf{vincolo}\index{vincolo!CP} \`e una relazione fra
una o pi\`u variabili che limita le combinazioni di valori
ammissibili. Esempi:
\begin{itemize}
\item \emph{Aritmetico}: $x + y \leq 10$.
\item \emph{Disuguaglianza}: $x \neq y$.
\item \emph{Globale}: \texttt{AllDifferent}$(x_1, \ldots, x_n)$
  significa ``i valori assunti da queste variabili devono
  essere tutti distinti''. Equivale a $\binom{n}{2}$ vincoli
  $x_i \neq x_j$, ma il propagatore globale \`e molto pi\`u
  efficiente (sfrutta il teorema di Hall, vedi
  cap.~\ref{ch:metodo_hall}).
\end{itemize}
\end{defbox}

\subsection{Propagazione}

\begin{defbox}
La \textbf{propagazione}\index{propagazione (CP)} \`e
l'applicazione di regole locali che, vista una scelta o
restrizione su una variabile, restringono il dominio di
\emph{altre} variabili senza che si debba fare alcun
``tentativo''. Si tratta di \emph{deduzione automatica}.
\end{defbox}

\begin{esempiobox}
Tre variabili $x, y, z \in \{1, 2, 3\}$ e tre vincoli:
$x \neq y$, $y \neq z$, $z \neq x$. Supponi di aver fissato
$x = 1$. La propagazione fa:
\begin{enumerate}
\item Da $x \neq y$: rimuove $1$ dal dominio di $y$. Ora
  $y \in \{2, 3\}$.
\item Da $z \neq x$: rimuove $1$ dal dominio di $z$. Ora
  $z \in \{2, 3\}$.
\end{enumerate}
A questo punto se aggiungi $y = 2$, propaga ancora:
\begin{enumerate}
\setcounter{enumi}{2}
\item Da $y \neq z$: rimuove $2$ dal dominio di $z$. Ora
  $z = 3$ (forzato).
\end{enumerate}
La propagazione ha dedotto $z$ senza enumerare. Senza
propagazione il solver avrebbe provato $z = 1, 2, 3$ uno per
uno e fallito su due dei tre tentativi.
\end{esempiobox}

\subsection{Ricerca}

\begin{defbox}
La \textbf{ricerca}\index{ricerca (CP)} \`e la fase in cui il
solver, esaurita la propagazione, sceglie a caso una
variabile e un valore, lo prova, e ricomincia a propagare. Se
la propagazione raggiunge una contraddizione (un dominio
vuoto), la ricerca \emph{fa backtracking}: torna indietro,
elimina quel valore dal dominio della variabile scelta, e
prova un altro.
\end{defbox}

La sequenza ``decisione $\to$ propagazione $\to$ eventuale
backtracking'' \`e il cuore del solver. La differenza fra un
solver lento e uno veloce sta in:
\begin{itemize}
\item \textbf{quale variabile scelgo}? Euristiche tipiche:
  ``la pi\`u vincolata'' (smallest domain first), ``la pi\`u
  attiva nei conflitti recenti'' (VSIDS).
\item \textbf{quale valore provo}? ``Il pi\`u promettente''
  (last solution first), o random.
\item \textbf{quanto imparo da ogni conflitto}? CDCL avanzato
  con clause learning + restart + clause forgetting.
\end{itemize}

\subsection{Combinazione: il ciclo CP-SAT}

\begin{figure}[h]
\centering
\begin{tikzpicture}[node distance=1.0cm and 0.8cm,
  block/.style={draw, very thick, rounded corners=4pt,
                minimum width=3.4cm, minimum height=0.9cm,
                align=center, font=\small},
  arr/.style={-{Stealth[length=2.5mm]}, thick}]
  \node[block, fill=cBox!40] (init) {Modello iniziale\\\footnotesize variabili + vincoli};
  \node[block, fill=cPref!15, below=of init] (prop) {Propagazione};
  \node[block, fill=cOro!20, below=of prop] (decide) {Decisione: scegli\\variabile e valore};
  \node[block, fill=cEnf!15, below left=1cm and -0.6cm of decide] (sol) {Soluzione\\trovata};
  \node[block, fill=cHard!15, below right=1cm and -0.6cm of decide] (back) {Conflitto:\\backtrack +\\impara clausola};
  \draw[arr] (init) -- (prop);
  \draw[arr] (prop) -- (decide);
  \draw[arr] (decide.south) -- node[draw=none, left, font=\scriptsize]
        {tutto assegnato} (sol);
  \draw[arr] (decide.south) -- node[draw=none, right, font=\scriptsize]
        {dominio vuoto} (back);
  \draw[arr] (back) to[bend right=45] (prop);
\end{tikzpicture}
\caption{Il ciclo CP-SAT. Dopo aver propagato il modello
iniziale, il solver alterna decisioni e propagazioni;
ogni conflitto produce una clausola appresa che evita di
ripetere lo stesso errore.}
\label{fig:cpsat-cycle}
\end{figure}

La fig.~\ref{fig:cpsat-cycle} mostra il ciclo. Importante:
\emph{ogni conflitto non \`e tempo perso}, perch\'e produce una
clausola appresa che diventa parte permanente del modello
e accelera tutte le ricerche future.

\section{Esempio mini-completo: 3 docenti, 2 classi}

Per \emph{vedere} CP-SAT all'opera, prendiamo un orario
miniatura: 3 docenti (Adami, Bruni, Conti), 2 classi (1A, 1B),
4 slot (luned\`i 1, luned\`i 2, marted\`i 1, marted\`i 2).
Materie e monte ore:
\begin{itemize}
\item 1A: matematica 2h, italiano 1h, storia 1h.
\item 1B: matematica 1h, italiano 2h, storia 1h.
\end{itemize}
Assegnazioni docenti--materia (gi\`a fatte da Phase A):
\begin{itemize}
\item Adami insegna matematica in 1A (2h) e in 1B (1h).
\item Bruni insegna italiano in 1A (1h) e in 1B (2h).
\item Conti insegna storia in 1A (1h) e in 1B (1h).
\end{itemize}
Vincolo aggiuntivo: Bruni non lavora luned\`i 1
(indisponibilit\`a HARD).

\textbf{Modellazione.} Per ogni (classe, materia, ora) ho una
variabile intera che dice in quale slot va. Per esempio
$\texttt{S}_{1\text{A,mate,1}} \in \{1,2,3,4\}$ (l'ora di
matematica della 1A in lunedi-1, lunedi-2, marted\`i-1 o
marted\`i-2). In totale ho 8 variabili: due ore di mate-1A,
una di ita-1A, una di sto-1A, una di mate-1B, due di ita-1B,
una di sto-1B.

Vincoli HARD:
\begin{enumerate}
\item Una classe non pu\`o avere due lezioni nello stesso
  slot: \texttt{AllDifferent} sui 4 slot di ogni classe.
\item Un docente non pu\`o tenere due lezioni nello stesso
  slot: per ogni coppia di lezioni dello stesso docente,
  $\texttt{S}_a \neq \texttt{S}_b$.
\item Bruni non lavora luned\`i 1: i due slot di ita-1B e
  l'uno di ita-1A non possono valere $1$.
\end{enumerate}

\textbf{Propagazione iniziale.} Dal vincolo (3): tre variabili
perdono subito $1$ dai loro domini. Ora ita-1A $\in
\{2,3,4\}$, e analogamente per le due ita-1B.

\textbf{Prima decisione.} Il solver sceglie la variabile con
dominio pi\`u piccolo. Tutte le 8 hanno dominio simile, ma le
3 di Bruni ne hanno solo 3. Sceglie $\texttt{ita-1A} = 2$
(luned\`i-2, primo valore disponibile).

\textbf{Propagazione dopo la decisione.} Il vincolo
\texttt{AllDifferent} sulla 1A rimuove $2$ dal dominio delle
altre 3 variabili della 1A. Il vincolo ``Bruni non in due
slot'' rimuove $2$ dalle due variabili di ita-1B. Ora ita-1B
ha dominio $\{3,4\}$ \emph{e devono essere entrambi distinti}:
in pratica $\{3,4\}$ \`e l'unica assegnazione possibile.

In altre parole, la propagazione ha gi\`a forzato i due slot
di Bruni in 1B. Resta da piazzare 5 variabili: mate-1A (2),
sto-1A (1), mate-1B (1), sto-1B (1).

\textbf{Seconda decisione.} Il solver tenta
$\texttt{mate-1A,1} = 1$. Propaga: rimuove $1$ da mate-1A,2
(dominio: $\{3,4\}$), da sto-1A (dominio $\{3,4\}$), e da
mate-1B (Adami non pu\`o essere in due posti, dominio
$\{2,3,4\}$, ma il $2$ \`e gi\`a stato preso da nessuna ora di
Adami quindi rimane $\{2,3,4\}$).

\textbf{Successive decisioni.} Il solver continua. Tipicamente
in 5--10 step di decisione+propagazione tutte le 8 variabili
sono fissate. Se in qualche step trovasse $\texttt{mate-1A,2}
= 3$ ma sto-1A $= 3$, il vincolo \texttt{AllDifferent}
darebbe contraddizione: backtracking.

In questo esempio il solver trova soluzione in pochi
millisecondi, senza neanche dover imparare clausole. Il punto
\`e che il \emph{numero di tentativi} \`e drasticamente ridotto
dalla propagazione: invece di esaminare $4^8 = 65{,}536$
combinazioni, il solver ne esamina meno di 20.

\begin{esempiobox}
\textbf{Una soluzione possibile}:
\begin{center}
\begin{tabular}{lcccc}
\toprule
& Lun-1 & Lun-2 & Mar-1 & Mar-2 \\
\midrule
1A & mate (Adami) & ita (Bruni) & sto (Conti) & mate (Adami) \\
1B & sto (Conti) & mate (Adami) & ita (Bruni) & ita (Bruni) \\
\bottomrule
\end{tabular}
\end{center}
Ogni classe ha 4 ore distinte; Adami non sta mai in due
posti contemporanei (lun-1, lun-2, mar-2 sono i suoi 3
slot); Bruni non insegna luned\`i 1; Conti insegna lun-1 e
mar-1. Tutto rispetta i vincoli.
\end{esempiobox}

\section{Generalizzazione alla scuola reale}

In una scuola di 30 classi e 60 docenti la stessa logica
funziona, con due differenze importanti:

\begin{enumerate}
\item \textbf{La propagazione fa gran parte del lavoro}. I
  vincoli \texttt{AllDifferent}, le indisponibilit\`a, le
  preferenze rigide e i lab condivisi creano una rete molto
  fitta di deduzioni: una scelta su una classe del triennio
  spesso costringe decine di altri slot.
\item \textbf{La ricerca diventa interessante quando i vincoli
  HARD sono ``quasi soddisfatti''}. Quando il solver ha
  piazzato il 70\% delle lezioni e si trova davanti alle
  ultime, qualunque scelta porta a contraddizioni e il
  backtracking si scatena. \`E qui che le clausole apprese
  fanno la differenza: dopo qualche centinaio di conflitti il
  solver ``capisce'' che certe combinazioni non funzionano e
  smette di provarle.
\end{enumerate}

\subsection{L'obiettivo: ottimizzazione, non solo
soddisfacibilit\`a}

CP-SAT non si limita a trovare \emph{una} soluzione: pu\`o
\emph{ottimizzare}. Si dichiara una variabile aggiuntiva
$\texttt{obj}$ definita come somma pesata di pezzi
``brutti'' (es. ore isolate dei docenti, sesta ora delle
classi, lezioni in slot non preferiti) e si chiede al solver
di minimizzarla. Internamente CP-SAT alterna ricerca della
soluzione e ricerca di una migliore: ogni volta che ne trova
una, aggiunge il vincolo $\texttt{obj} < \text{best}$ e
ricomincia. Termina quando dimostra che non esiste soluzione
migliore o quando finisce il tempo.

In $\pi$Tantum la funzione obiettivo \`e descritta nel
cap.~\ref{ch:workflow_di_ottimizzazione}; si veda anche la
documentazione DSL per la sintassi degli obiettivi
componibili.

\section{Quando usarlo (e quando no)}

\textbf{CP-SAT brilla quando}:
\begin{itemize}
\item Il problema \`e \emph{combinatorio puro} (variabili
  intere, vincoli logici): scheduling, configurazione,
  packing, routing piccoli/medi.
\item I vincoli sono \emph{molto interconnessi}: la
  propagazione ha materiale su cui lavorare.
\item Il numero di variabili \`e nell'ordine $10^3$--$10^5$:
  oltre, conviene decomporre o usare metaeuristiche.
\end{itemize}

\textbf{Limiti}:
\begin{itemize}
\item Sopra $\sim 10^6$ variabili boolean equivalenti, anche
  CP-SAT diventa lento. \`E uno dei motivi per cui
  $\pi$Tantum decompone le scuole grandi (cap.~\ref{ch:metodo_spettrale}).
\item Sui problemi \emph{continui} (variabili reali) CP-SAT
  non lavora; servono solver LP/MIP/SOCP.
\item Se l'obiettivo \`e una somma di molti termini con pesi
  fini, la ricerca dell'ottimo esatto pu\`o impiegare a lungo:
  spesso conviene fermarla con un \texttt{time\_limit} e
  raffinare con metaeuristiche.
\end{itemize}

\begin{avvertenzabox}
CP-SAT \`e una ``scatola intelligente'' ma non magica.
Modellare male un problema (es. con vincoli ridondanti, o con
variabili intere che potevano essere booleane) lo rallenta di
ordini di grandezza. La regola pratica \`e: \emph{usa i
vincoli globali quando puoi}, perch\'e ogni vincolo globale ha
un propagatore specializzato che lavora molto meglio di
$n^2$ vincoli binari.
\end{avvertenzabox}

\section{Connessioni con il resto del sistema}

\begin{connessionebox}
\begin{itemize}
\item \textbf{Decomposizione spettrale}
  (cap.~\ref{ch:metodo_spettrale}): per scuole oltre $\sim 80$
  classi, $\pi$Tantum spezza il problema in cluster e
  invoca CP-SAT su ognuno.
\item \textbf{LNS / SA / TS / ILS} (cap.~\ref{ch:metodo_metaeuristiche}):
  ogni metaeuristica usa CP-SAT per ``riparare''
  sotto-vicinati ad ogni iterazione (ALNS in particolare).
\item \textbf{Hall pre-check} (cap.~\ref{ch:metodo_hall}):
  prima di lanciare CP-SAT, $\pi$Tantum verifica che
  l'assegnamento docenti $\to$ classi sia \emph{matching-feasible}.
  Se Hall fallisce, CP-SAT non potrebbe trovare nessuna
  soluzione e darebbe \texttt{INFEASIBLE} dopo aver bruciato
  minuti di calcolo.
\item \textbf{Lagrangian Relaxation}
  (cap.~\ref{ch:metodo_lagrangian}): rilassa alcuni vincoli
  HARD trasformandoli in penalit\`a sull'obiettivo, riducendo
  il problema CP-SAT a sotto-problemi pi\`u piccoli.
\end{itemize}
\end{connessionebox}

\begin{biblioparvabox}
\begin{itemize}
\item Lettura introduttiva (divulgativa): Russell \& Norvig,
  \emph{Artificial Intelligence: A Modern Approach},
  capp.~6--7 sui CSP.
\item Manuale di riferimento: Rossi, Van Beek, Walsh,
  \emph{Handbook of Constraint Programming}~\cite{rossi2006},
  capp.~3 (CSP), 4 (consistency), 6 (search).
\item Storia di CDCL e SAT moderno: Marques-Silva, Lynce,
  Malik, \emph{Conflict-Driven Clause Learning SAT
  Solvers}, in \emph{Handbook of Satisfiability}, IOS Press
  2009.
\item Documentazione OR-Tools: \url{https://developers.google.com/optimization/cp/cp_solver}
  -- molti esempi pratici in Python.
\item Contesto teorico (NP-completezza):
  Garey \& Johnson~\cite{garey1979} resta il classico.
\end{itemize}
\end{biblioparvabox}
